\section*{Capítulo \ref{ch:g8:combustion-and-fuels} --- La combustión: combustibles, productos y gases de efecto invernadero}

\begin{solution}{exo:g8:combustion-and-fuels:1}
\ce{C + O2 -> CO2}.
\end{solution}

\begin{solution}{exo:g8:combustion-and-fuels:2}
Con oxígeno suficiente, todo el carbono se convierte en dióxido de
carbono (completa). Con demasiado poco, una parte se convierte en
monóxido de carbono o en hollín (incompleta).
\end{solution}

\begin{solution}{exo:g8:combustion-and-fuels:3}
No tiene color, ni olor, ni sabor: nadie lo nota y, sin embargo, es
tóxico.
\end{solution}

\begin{solution}{exo:g8:combustion-and-fuels:4}
Fósiles: el carbón, el gas natural, la gasolina. Renovables: la madera,
el biogás, el etanol de remolacha.
\end{solution}

\begin{solution}{exo:g8:combustion-and-fuels:5}
\ce{C3H8 + 5O2 -> 3CO2 + 4H2O}.
\end{solution}

\begin{solution}{exo:g8:combustion-and-fuels:6}
Carbono: 2 \ce{CO2}; hidrógeno: 3 \ce{H2O}; oxígeno a la derecha, $4 + 3 =
7$, de los que 1 viene del etanol, así que 3 \ce{O2}:
\ce{C2H6O + 3O2 -> 2CO2 + 3H2O}.
\end{solution}

\begin{solution}{exo:g8:combustion-and-fuels:7}
Hollín, es decir, carbono. La \omterm{def:g4:what-a-fire-needs:burning}{combustión} es incompleta: a la llama le
falta oxígeno.
\end{solution}

\begin{solution}{exo:g8:combustion-and-fuels:8}
Unas 339 ppm en 1980 y unas 370 ppm en 2000: un aumento de unas 30 ppm.
\end{solution}

\begin{solution}{exo:g8:combustion-and-fuels:9}
El carbono de la madera lo tomó del \omterm{def:g4:air-a-mixture-of-gases:air}{aire} el árbol hace unos años o unas
décadas: al quemarla se devuelve ese dióxido de carbono. El carbono del
carbón estuvo almacenado bajo tierra durante millones de años: al
quemarlo se añade dióxido de carbono nuevo al \omterm{def:g4:air-a-mixture-of-gases:air}{aire}.
\end{solution}

\begin{solution}{exo:g8:combustion-and-fuels:10}
\ce{2C + O2 -> 2CO}.
\end{solution}

\begin{solution}{exo:g8:combustion-and-fuels:11}
$(427 - 316)/316 = 111/316 \approx 0.35$: alrededor del \qty{35}{\percent}.
\end{solution}

\begin{solution}{exo:g8:combustion-and-fuels:12}
Al principio hay oxígeno suficiente: \ce{CH4 + 2O2 -> CO2 + 2H2O}. A
medida que se gasta el oxígeno de la habitación cerrada, la llama \omterm{def:g4:what-a-fire-needs:burning}{arde}
con poco \omterm{def:g4:air-a-mixture-of-gases:air}{aire}:
\ce{2CH4 + 3O2 -> 2CO + 4H2O}. El monóxido de carbono, inodoro y
tóxico, se acumula en la habitación.
\end{solution}

\begin{solution}{pb:g8:combustion-and-fuels:1}
\textbf{1.} \ce{C3H8 + 5O2 -> 3CO2 + 4H2O}.

\textbf{2.} Izquierda: 3 C, 8 H, 10 O. Derecha: 3 C, $4 \times 2 = 8$ H,
$3 \times 2 + 4 = 10$ O. Ajustada.

\textbf{3.} \ce{2C3H8 + 7O2 -> 6CO + 8H2O}.

\textbf{4.} La \omterm{def:g8:combustion-and-fuels:complete}{combustión completa}: 5 \omterm{def:g7:atoms-and-molecules:molecule}{moléculas} de oxígeno por \omterm{def:g7:atoms-and-molecules:molecule}{molécula}
de propano, frente a $7/2 = 3.5$ en la incompleta.

\textbf{5.} El \omterm{def:g4:air-a-mixture-of-gases:air}{aire} de la tienda cerrada pronto se queda sin oxígeno; a
la llama ya no le llega suficiente y \omterm{def:g4:what-a-fire-needs:burning}{arde} de forma incompleta, amarilla.

\textbf{6.} El monóxido de carbono: no tiene color, ni olor (ni sabor).

\textbf{7.} GHS02, la llama: extremadamente inflamable. GHS04, la
bombona de gas: gas guardado a presión.

\textbf{8.} Por ejemplo: usarlo solo al \omterm{def:g4:air-a-mixture-of-gases:air}{aire} libre, nunca en una tienda
de campaña ni en una habitación cerrada; mantenerlo lejos de todo lo que
pueda \omterm{def:g4:what-a-fire-needs:burning}{arder}; cambiar el cartucho lejos de cualquier llama.

\textbf{9.} $132 \div 44 = 3$ veces más pesado.

\textbf{10.} Del oxígeno del \omterm{def:g4:air-a-mixture-of-gases:air}{aire}, que acaba en el dióxido de carbono (y
en el agua).

\textbf{11.} Un \omterm{def:g8:combustion-and-fuels:fossil}{combustible fósil} (procede del gas natural o del
petróleo): al quemarlo se añade dióxido de carbono al \omterm{def:g4:air-a-mixture-of-gases:air}{aire}.

\textbf{12.} $450 \times 3 = \qty{1350}{g} = \qty{1.35}{kg}$ de dióxido
de carbono.
\end{solution}
